\documentclass[UTF8,11pt,a4paper]{ctexart}

\usepackage[margin=1.9cm,headheight=15pt]{geometry}
\usepackage{amsmath,amssymb}
\usepackage{booktabs}
\usepackage{longtable}
\usepackage{array}
\usepackage{tabularx}
\usepackage{enumitem}
\usepackage{xcolor}
\usepackage{tcolorbox}
\usepackage{fancyhdr}
\usepackage{hyperref}
\hypersetup{colorlinks=true,linkcolor=blue!45!black,urlcolor=blue!45!black}

\definecolor{ink}{HTML}{1F2937}
\definecolor{accent}{HTML}{B45309}
\definecolor{muted}{HTML}{6B7280}

\ctexset{
  section    = {format = {\large\bfseries\color{accent}}},
  subsection = {format = {\normalsize\bfseries\color{ink}}},
}

\pagestyle{fancy}
\fancyhf{}
\fancyhead[L]{\small\color{muted}SOE104 CW2 · 操作手册}
\fancyhead[R]{\small\color{muted}IntelliJ IDEA 配置}
\fancyfoot[C]{\small\color{muted}\thepage}
\renewcommand{\headrulewidth}{0.4pt}

\setlist{itemsep=1pt,topsep=2pt,parsep=0pt}
\setlength{\parskip}{0.28em}
\renewcommand{\arraystretch}{1.2}

\newtcolorbox{stepbox}[1][]{
  colback=blue!2, colframe=blue!45!black, boxrule=0.6pt, arc=2pt,
  left=6pt, right=6pt, top=4pt, bottom=4pt, fonttitle=\bfseries, #1
}
\newtcolorbox{keybox}[1][]{
  colback=orange!4, colframe=orange!55!black, boxrule=0.6pt, arc=2pt,
  left=6pt, right=6pt, top=4pt, bottom=4pt, fonttitle=\bfseries, #1
}
\newtcolorbox{warnbox}[1][]{
  colback=red!3, colframe=red!45!black, boxrule=0.6pt, arc=2pt,
  left=6pt, right=6pt, top=4pt, bottom=4pt, fonttitle=\bfseries, #1
}
\newtcolorbox{goodbox}[1][]{
  colback=green!3, colframe=green!35!black, boxrule=0.6pt, arc=2pt,
  left=6pt, right=6pt, top=4pt, bottom=4pt, fonttitle=\bfseries, #1
}

\newcommand{\code}[1]{\texttt{#1}}
\newcommand{\menu}[1]{\textsf{#1}}

\begin{document}

\begin{center}
{\LARGE\bfseries IntelliJ IDEA 配置手册}\\[4pt]
{\large 打开项目 · 导入 jar · 验证通过}\\[6pt]
{\small\color{muted}SOE104 CW2 \quad 2026-11-21 \quad 整理：A}
\end{center}

\vspace{-0.2cm}

\begin{keybox}[title={预期结果}]
做完这本手册，你应该看到：\textbf{打开 \code{.java} 文件时没有红色波浪线}，
\code{import edu.soe104.oregon.api.*;} 能自动补全，
\menu{Build} $\to$ \menu{Build Project} 能成功。

\textbf{注意：你不能在本地运行这个游戏。}jar 里只有接口定义，没有引擎。
本地只负责“编译通过”，跑游戏只能在平台上。这件事后面还会再强调一次。
\end{keybox}

%=====================================================================
\section{第一步：用 Open，不要用 New Project}

\begin{warnbox}[title={这是新手最常犯的错}]
\textbf{不要点 \menu{New Project}}。那是从零建一个空项目，
会覆盖掉仓库里已有的目录结构。

我们要做的是\textbf{打开一个已经存在的文件夹}。
\end{warnbox}

\begin{stepbox}[title={步骤}]
\begin{enumerate}[leftmargin=1.8em]
  \item 先把仓库克隆到本地（见 \texttt{06\_git\_handbook.pdf} 第二节）。
        假设路径是 \code{D:\textbackslash XJTLU\textbackslash 2026 Y2 S1\textbackslash SOE104\textbackslash CW2\textbackslash soe104-cw2}
  \item 打开 IDEA，在欢迎界面点 \menu{Open}（不是 New Project）。
  \item 选中 \code{soe104-cw2} 这个文件夹，点 \menu{OK}。
  \item 如果弹出 “Trust and Open Project?”，点 \menu{Trust Project}。
\end{enumerate}
\end{stepbox}

\begin{goodbox}[title={成功的标志}]
左侧项目树能看到 \code{src}、\code{lib}、\code{docs} 三个文件夹。
\end{goodbox}

%=====================================================================
\section{第二步：设置 JDK 21}

\noindent 课程用 JDK 21。你电脑上可能装了别的版本，要先指定。

\begin{stepbox}[title={步骤}]
\begin{enumerate}[leftmargin=1.8em]
  \item \menu{File} $\to$ \menu{Project Structure}（快捷键 \code{Ctrl+Alt+Shift+S}）。
  \item 左侧选 \menu{Project}。
  \item \menu{SDK} 下拉框：如果有 \code{21}，直接选它，跳到下一步。
  \item 如果没有，点 \menu{Add SDK} $\to$ \menu{Download JDK}，
        版本选 \code{21}，Vendor 选 \code{Oracle OpenJDK} 或 \code{Eclipse Temurin}，点 \menu{Download}。
  \item \menu{Language level} 设为 \code{21}。
  \item 点 \menu{Apply} $\to$ \menu{OK}。
\end{enumerate}
\end{stepbox}

\begin{warnbox}[title={版本不对会怎样}]
如果 Language level 低于 21，有些语法会报错；
如果高于 21，本地能编过但平台是 JDK 21，可能出现“本地能跑、平台报错”。
\textbf{统一用 21。}
\end{warnbox}

%=====================================================================
\section{第三步：把 src 标记为源码目录}

\noindent 告诉 IDEA“这个文件夹里是源代码”，否则它不会帮你检查语法。

\begin{stepbox}[title={步骤}]
\begin{enumerate}[leftmargin=1.8em]
  \item 在左侧项目树里，\textbf{右键}点 \code{src} 文件夹。
  \item 选 \menu{Mark Directory as} $\to$ \menu{Sources Root}。
  \item 成功后 \code{src} 图标会变成\textbf{蓝色}。
\end{enumerate}
\end{stepbox}

%=====================================================================
\section{第四步：导入 oregon-api.jar（关键步骤）}

\noindent 这一步做完，红色波浪线才会消失。

\begin{stepbox}[title={步骤}]
\begin{enumerate}[leftmargin=1.8em]
  \item \menu{File} $\to$ \menu{Project Structure}（\code{Ctrl+Alt+Shift+S}）。
  \item 左侧选 \menu{Libraries}。
  \item 点中间的 \menu{+} 号 $\to$ 选 \menu{Java}。
  \item 在弹出的文件选择框里，找到 \code{soe104-cw2\textbackslash lib\textbackslash oregon-api.jar}，
        选中它，点 \menu{OK}。
  \item 弹窗问 “Add to module?”，点 \menu{OK}（默认选中当前模块就行）。
  \item 点 \menu{Apply} $\to$ \menu{OK}。
\end{enumerate}
\end{stepbox}

\begin{goodbox}[title={成功的标志}]
在 \menu{Libraries} 列表里能看到 \code{oregon-api}，展开后有一堆 \code{.class} 文件。
\end{goodbox}

\begin{warnbox}[title={如果项目树里看不到 lib 文件夹}]
说明 jar 没被放进仓库。让 A 确认 \code{lib/oregon-api.jar} 已经提交，
然后你 \code{git pull} 一次。
\end{warnbox}

%=====================================================================
\section{第五步：设置 UTF-8 编码}

\noindent 不设的话，中文注释会变乱码，提交到平台可能编译失败。

\begin{stepbox}[title={步骤}]
\begin{enumerate}[leftmargin=1.8em]
  \item \menu{File} $\to$ \menu{Settings}（\code{Ctrl+Alt+S}）。
  \item 左侧展开 \menu{Editor} $\to$ \menu{File Encodings}。
  \item 三个下拉框全部设为 \code{UTF-8}：
        \begin{itemize}[leftmargin=1.2em]
          \item \menu{Global Encoding}
          \item \menu{Project Encoding}
          \item \menu{Default encoding for properties files}
        \end{itemize}
  \item 勾选 \menu{Transparent native-to-ascii conversion}（可选，但建议勾）。
  \item 点 \menu{Apply} $\to$ \menu{OK}。
\end{enumerate}
\end{stepbox}

%=====================================================================
\section{第六步：验证（做完这一步才算成功）}

\subsection{方法一：看有没有红线}

\begin{stepbox}[title={步骤}]
\begin{enumerate}[leftmargin=1.8em]
  \item 展开 \code{src}，双击任意一个 \code{.java} 文件。
  \item 看第一行 \code{import edu.soe104.oregon.api.*;} —— \textbf{应该是灰色的，没有红色波浪线}。
  \item 把光标放到 \code{Member} 上按 \code{Ctrl+B}，能跳转到 \code{.class} 文件就说明 jar 生效了。
\end{enumerate}
\end{stepbox}

\subsection{方法二：编译一次}

\begin{stepbox}[title={步骤}]
\begin{enumerate}[leftmargin=1.8em]
  \item 点菜单 \menu{Build} $\to$ \menu{Build Project}。
  \item 看底部 \menu{Build} 窗口，应该显示 \code{Build completed successfully}。
\end{enumerate}
\end{stepbox}

\begin{center}
\begin{tabular}{@{}p{6.6cm}p{6.6cm}@{}}
\toprule
\textbf{你看到的} & \textbf{说明} \\
\midrule
\code{Build completed successfully} & \textbf{成功}，可以开始写代码了 \\
\code{cannot find symbol} & jar 没导进来，回到第四步 \\
\code{package ... does not exist} & 同上 \\
\code{invalid source release: 21} & JDK 版本不对，回到第二步 \\
\bottomrule
\end{tabular}
\end{center}

%=====================================================================
\section{第七步：新建类文件（注意不要建包）}

\begin{warnbox}[title={绝对不要创建 package}]
课程的提交要求是\textbf{没有 package 声明的默认包}。
如果你在 IDEA 里新建类时填了包名（比如 \code{com.team}），
文件第一行会出现 \code{package com.team;}——\textbf{提交会直接编译失败}。
\end{warnbox}

\begin{stepbox}[title={步骤}]
\begin{enumerate}[leftmargin=1.8em]
  \item 右键点 \code{src} 文件夹 $\to$ \menu{New} $\to$ \menu{Java Class}。
  \item 在弹出的框里\textbf{只填类名}，例如 \code{TrailHunter}。
  \item \textbf{不要}在 \code{TrailHunter} 前面加任何路径或包名。
  \item 点 \menu{Enter}。文件应该在 \code{src} 根目录，而不是某个子文件夹里。
\end{enumerate}
\end{stepbox}

\begin{goodbox}[title={正确与错误对比}]
\textbf{正确：}文件在 \code{src/TrailHunter.java}，第一行是 \code{import ...}\\
\textbf{错误：}文件在 \code{src/com/team/TrailHunter.java}，第一行是 \code{package com.team;}
\end{goodbox}

\begin{keybox}[title={IDEA 的一个提示}]
新建默认包里的类时，IDEA 可能提示 \menu{Warning: class is in the default package}。
\textbf{这个警告可以忽略}——课程本来就要求默认包。
\end{keybox}

%=====================================================================
\section{常见问题}

\begin{center}
\small
\begin{longtable}{@{}p{5.2cm}p{7.8cm}@{}}
\toprule
\textbf{现象} & \textbf{原因与解决} \\
\midrule
\endhead
所有 import 都是红的 & jar 没加。回到第四步 \\
\code{src} 图标是灰色 & 没标记 Sources Root。回到第三步 \\
中文注释显示成方框 & 编码不是 UTF-8。回到第五步 \\
Build 报 \code{invalid source release} & JDK 不是 21。回到第二步 \\
新建的类跑到子文件夹里 & 建的时候填了包名。删掉重建 \\
项目树里没有 \code{lib} & 仓库没拉全，\code{git pull} 一次 \\
找不到 \menu{Project Structure} & 快捷键 \code{Ctrl+Alt+Shift+S}，或 \menu{File} 菜单里找 \\
\bottomrule
\end{longtable}
\end{center}

%=====================================================================
\section{三件不能做的事}

\begin{warnbox}[title={记住这三条}]
\begin{enumerate}[leftmargin=1.8em]
  \item \textbf{不要在本地“运行”游戏。}jar 里没有引擎，\code{main} 方法也没有。
        点运行只会报错，这不是你的问题。跑游戏只能在平台上。
  \item \textbf{不要新建 package。}提交必须是无包声明的默认包。
  \item \textbf{不要改 \code{lib/oregon-api.jar}。}它只读，也不需要改。
\end{enumerate}
\end{warnbox}

\begin{keybox}[title={本地和平台的职责划分}]
\begin{center}
\begin{tabular}{@{}p{3.0cm}p{4.6cm}p{4.6cm}@{}}
\toprule
& \textbf{本地（IDEA）} & \textbf{平台} \\
\midrule
语法检查 & \checkmark & — \\
自动补全 & \checkmark & — \\
能否编译 & 近似 & 权威 \\
跑游戏 & \textbf{不能} & \checkmark \\
禁用 API 扫描 & 不能 & \checkmark \\
打分 & 不能 & \checkmark \\
\bottomrule
\end{tabular}
\end{center}
\end{keybox}

\vfill
\begin{center}
\small\color{muted}
配套手册：\texttt{06\_git\_handbook.pdf}（Git 一条龙）
\end{center}

\end{document}
